\thispagestyle{empty}
\vspace*{18mm}
\begin{center}
  {\sffamily\small ENGINEERING \& DATA SCIENCES \,$\cdot$\, CHEMISTRY S1-1}\\[10pt]
  {\sffamily\bfseries\Huge\color{heading} From atomic entities}\\[4pt]
  {\sffamily\bfseries\Huge\color{heading} to physico-chemical systems}\\[12pt]
  {\sffamily\LARGE Study guide}\\[4pt]
  {\large Concepts in detail \,$\cdot$\, Recaps \,$\cdot$\, Worksheets}\\[3pt]
  {\normalsize Semester 1 \,$\cdot$\, Academic Year 2026--2027}
\end{center}
\vspace{10mm}
{\sffamily\bfseries\color{heading} What is in this guide}
\begin{itemize}[leftmargin=7mm, itemsep=3pt]
  \item Five units, one per chapter of the student handout (pages 1--20). Each unit has three parts:
  \begin{itemize}
    \item \textbf{Concepts in detail}: every notion of the chapter with a precise definition, the
          reasoning behind it, worked examples and common mistakes.
          \colorbox{fillcol!20}{Orange boxes} complete the blank boxes and tables of the handout.
    \item \textbf{Recap}: one page with the definitions, formulas, key numbers and a checklist.
    \item \textbf{Worksheet}: 15 questions that get progressively harder.
  \end{itemize}
  \item Worked answers to every question are in a separate booklet,
        \emph{Chem-S1-1\_Worksheet-Answers.pdf}. Try the questions first.
\end{itemize}
\vspace{4mm}
{\sffamily\bfseries\color{heading} Contents}\par\vspace{2mm}
\begin{center}
\renewcommand{\arraystretch}{1.3}
\begin{tabular}{@{}c l c@{}}
\toprule
Unit & Title & Handout\\ \midrule
1 & Classical description of the atom & ch.\ 1 (1.1--1.3)\\
2 & Quantum model of the atom & ch.\ 2 (2.1--2.5)\\
3 & Quantum numbers and atomic orbitals & ch.\ 3 (3.1--3.4)\\
4 & Electronic configuration and periodic table & ch.\ 4 (4.1--4.2), pp.\ 13--15\\
5 & Atomic properties and periodicity & ch.\ 5 (5.1--5.5), pp.\ 16--20\\ \bottomrule
\end{tabular}
\end{center}
\vspace{4mm}
{\sffamily\bfseries\color{heading} Worksheet levels}\par\vspace{1mm}
\begin{tabular}{@{}l l@{}}
\level{1} & Q1--4: recall a definition, apply one formula.\\
\level{2} & Q5--8: explain a result or chain two steps.\\
\level{3} & Q9--12: multi-step problems, interpreting data.\\
\level{4} & Q13--15: synthesis, links between units, going beyond the handout.
\end{tabular}
\vspace{4mm}
\begin{tcolorbox}[enhanced, colback=black!3, colframe=black!45, boxrule=0.4pt, arc=1mm, fontupper=\small]
Constants as in the handout: $c = 3.00\times10^8\ \mathrm{m\,s^{-1}}$, $h = 6.62\times10^{-34}$~J\,s,
$e = 1.6\times10^{-19}$~C, $R_H = 3.29\times10^{15}$~Hz, $a_0 = 52.9$~pm, $1\ \mathrm{amu} =
1.6605\times10^{-27}$~kg, $N_A = 6.022\times10^{23}\ \mathrm{mol^{-1}}$. Where the handout has a typo (for
example the Balmer formula, or ``650'' written for 700~nm on the spectra), the guide says so.
\end{tcolorbox}
\clearpage
